% Data Processing Utility
\documentclass{article}
\usepackage{array}
\usepackage{booktabs}
\usepackage{xparse}

\begin{document}

\title{Data Processing Utility}
\author{TEX Engine}
\date{}\maketitle

\section{Introduction}
This is a data processing utility written in \TeX. It demonstrates the processing of data structures within a LaTeX document.

\section{Data Processing}
We define a simple data structure to hold our records.
\begin{definition}[Data Record]
    A \textbf{Data Record} is a tuple of values:
    \begin{equation}
        R = (id, name, value)
    \end{equation}
\end{definition}

\section{Example Output}
We can process a sample dataset and format it as a table.

\begin{tabular}{lll}
\toprule
\textbf{ID} & \textbf{Name} & \textbf{Value} \\
\midrule
1 & Alpha & 100 \\
2 & Beta & 200 \\
3 & Gamma & 300 \\
\bottomrule
\end{tabular}

\section{Conclusion}
This utility provides a foundation for processing and presenting data using \TeX.

\end{document}
